\chapter*{Abstract}
\addcontentsline{toc}{chapter}{Abstract}

Physical AI systems do not act on reality directly. They act on measurements, filtered
state estimates, delayed telemetry, and learned summaries of the world. When that
observation surface degrades, an action can remain admissible on the observed state while
already being unsafe on the true state. This dissertation treats that discrepancy as an
action-semantics safety problem rather than as a prediction problem alone, and names it the
observation--action safety gap (OASG).

Existing lines of work address important parts of this failure mode: conformal uncertainty
quantifies predictive spread, safe control constrains actions, runtime assurance supervises
release, and fault monitoring detects degraded telemetry. What is still missing is the
integrated runtime object that connects those pieces at the moment of physical release.
ORIUS, Observation--Reality Integrity for Universal Safety, is proposed as that missing
object: a certificate-aware runtime safety layer for physical AI under degraded
observation.

The dissertation develops ORIUS as a disciplined two-tier progression. First, Chapters~1
through~12 establish the base three-domain CPU \texttt{DC3S} and \texttt{CertOS} runtime
program across \texttt{DomainAdapter} safety predicates, reliability-conditioned conformal
uncertainty, action repair, fallback semantics, and cryptographic certificate lifecycle
governance. On this foundational tier, the quantitative safety guarantees are stated as
finite-sample \emph{marginal} and reliability-stratified coverage bounds (\texttt{T1}--\texttt{T11})
validated on Battery Energy Storage (deep witness row), Autonomous Vehicles (\texttt{nuPlan}
runtime replay), and Medical ICU Monitoring (\texttt{MIMIC}/\texttt{BIDMC}).

Second, Chapter~13 introduces the \textbf{ORIUS + K-Bound} kinetic safety extension and its
Lean~4 mechanized proof suite (\texttt{Theorems~A, E, F, G, H} in
\texttt{formal/Orius/TheoremsFGH.lean}, zero \texttt{sorry}) to close four structural
stress boundaries: (i) unbounded telemetry blackouts via the Terminal Viability Safe-Stop
Kernel (\texttt{TV-SSK}, Theorem~F), which drives physical and financial states into an
absorbing zero-kinetic-energy safe equilibrium $\mathcal{Z}_{\mathrm{stop}}$ before residual
tube expiration; (ii) non-stationary and adversarial distribution shift via an
Anytime-Valid Ville Betting Martingale (\texttt{AnytimeEProcessGate}, Theorem~G) paired
with a breakdown-point-optimal Median-of-Means Byzantine sensor filter
(\texttt{ByzantineResilientOQE}); (iii) cyber-physical financial L2 flash-crash liquidation
via the Almgren--Chriss safe-liquidation isomorphism under an explicit liquidity floor
$L_{\mathrm{depth}}(w_t) \ge L_{\min} > 0$ (Theorem~H); and (iv) sub-microsecond C++17
inner-loop execution ($0.0056\,\mu\text{s}$ per step), 100,000-agent GPU swarm governance,
seven-domain real-data validation ($224,000$ steps, $0.000\%$ ORIUS TSVR), and closed-loop
mobile robot and electro-thermal battery bench verification. Chapter~14 synthesizes the
governance, reproducibility, and claim-locked boundary between the base three-domain CPU
release lane and the seven-domain \texttt{K-Bound} kinetic closure.
